\documentclass[aspectratio=169]{beamer}
\usepackage[utf8]{inputenc}
\usepackage[spanish]{babel}
\usepackage[T1]{fontenc}
\usepackage{lmodern}
\usepackage{tikz}

% ---------- Paleta verde ----------
\definecolor{verde1}{HTML}{00D9A5}
\definecolor{verde2}{HTML}{4DE8C0}
\definecolor{verde3}{HTML}{80F0D4}
\definecolor{verde4}{HTML}{B3F7E5}
\definecolor{verde5}{HTML}{E6FFF7}
\definecolor{verdeosc}{HTML}{00876A} % para texto legible

% ---------- Estilo general ----------
\usetheme{default}
\usefonttheme{professionalfonts}
\setbeamertemplate{navigation symbols}{}
\setbeamercolor{frametitle}{fg=black}
\setbeamerfont{frametitle}{series=\bfseries,size=\Large}
\setbeamercolor{structure}{fg=verde1}
\setbeamercolor{block title}{bg=verde1,fg=white}
\setbeamercolor{block body}{bg=verde5,fg=black}
\setbeamercolor{block title alerted}{bg=verdeosc,fg=white}
\setbeamercolor{block body alerted}{bg=verde4,fg=black}
\setbeamertemplate{blocks}[rounded][shadow=false]
\setbeamertemplate{itemize items}[circle]
\setbeamercolor{section in toc}{fg=black}

% Barra superior y pie con logo en diapositivas normales
\setbeamertemplate{frametitle}{%
  \vspace*{0.4cm}\insertframetitle\par
  \vspace*{-0.15cm}{\color{verde1}\rule{2.5cm}{2pt}}}
\setbeamertemplate{background}{%
  \begin{tikzpicture}[remember picture,overlay]
    \fill[verde1] (current page.north west) rectangle ([yshift=-0.15cm]current page.north east);
    \fill[verde2] ([xshift=-2.2cm]current page.south east) -- (current page.south east) -- ([yshift=1.6cm]current page.south east) -- cycle;
    \fill[verde4] ([xshift=-3.4cm]current page.south east) -- ([xshift=-2.2cm]current page.south east) -- ([yshift=1.6cm]current page.south east) -- ([yshift=2.2cm]current page.south east) -- cycle;
  \end{tikzpicture}}
\setbeamertemplate{footline}{%
  \hspace{0.4cm}\raisebox{0.15cm}{\includegraphics[height=0.55cm]{imagenes/logo_utec.png}}%
  \hfill{\scriptsize\color{gray}\insertframenumber}\hspace{1.2cm}\vspace{0.1cm}}

% ---------- Datos ----------
\newcommand{\etiq}[2]{{\color{verde1}\textbf{#1:}} \textbf{#2}\\[0.12cm]}

\begin{document}

% ================= CARÁTULA =================
{
\setbeamertemplate{background}{}
\setbeamertemplate{footline}{}
\begin{frame}[plain]
\begin{tikzpicture}[remember picture,overlay]
  % bloques diagonales verdes (lado derecho)
  \fill[verde5] ([xshift=6.0cm]current page.north) -- (current page.north east) -- (current page.south east) -- ([xshift=1.0cm]current page.south) -- cycle;
  \fill[verde4] ([xshift=7.2cm]current page.north) -- (current page.north east) -- (current page.south east) -- ([xshift=2.6cm]current page.south) -- ([xshift=4.6cm]current page.center) -- cycle;
  \fill[verde3] ([xshift=1.5cm]current page.center) -- ([xshift=4.6cm]current page.center) -- ([xshift=2.6cm]current page.south) -- ([xshift=-0.6cm]current page.south) -- cycle;
  \fill[verde2] ([xshift=-1.6cm]current page.north east) -- (current page.north east) -- ([yshift=-3.8cm]current page.north east) -- cycle;
  \fill[verde1] ([xshift=-0.9cm,yshift=-1.0cm]current page.north east) -- ([yshift=-1.0cm]current page.north east) -- ([yshift=-5.4cm]current page.north east) -- ([xshift=-0.9cm,yshift=-3.4cm]current page.north east) -- cycle;
  \fill[verde1] ([xshift=4.6cm]current page.center) -- ([xshift=6.8cm,yshift=1.2cm]current page.center) -- ([xshift=4.8cm,yshift=-4.0cm]current page.center) -- ([xshift=2.6cm]current page.south) -- cycle;

  % título
  \node[anchor=north west, text width=7cm, align=left, font=\LARGE\bfseries]
    at ([xshift=0.9cm,yshift=-0.9cm]current page.north west)
    {TÍTULO DE LA PRESENTACIÓN};

  % datos
  \node[anchor=north west, text width=8.5cm, font=\footnotesize, align=left]
    at ([xshift=0.9cm,yshift=-3.6cm]current page.north west) {%
    \etiq{CURSO}{NOMBRE DEL CURSO}
    \etiq{ASESOR}{NOMBRE DEL ASESOR}
    \etiq{ALUMNO}{NOMBRE DEL ALUMNO}
    \etiq{CICLO}{AAAA-N}};

  % logo
  \node[anchor=south west] at ([xshift=0.6cm,yshift=0.4cm]current page.south west)
    {\includegraphics[height=1.4cm]{imagenes/logo_utec.png}};
\end{tikzpicture}
\end{frame}
}

% ================= CONTENIDO =================
\begin{frame}{Contenido}
  \tableofcontents
\end{frame}

\section{Introducción}
\begin{frame}{Introducción}
  \begin{itemize}
    \item Primer punto de la introducción.
    \item Segundo punto.
    \item Tercer punto.
  \end{itemize}
\end{frame}

\section{Desarrollo}
\begin{frame}{Desarrollo}
  \begin{block}{Objetivo}
    Texto del objetivo.
  \end{block}
  \begin{alertblock}{Nota}
    Texto de la nota.
  \end{alertblock}
\end{frame}

\begin{frame}{Dos columnas}
  \begin{columns}
    \column{0.5\textwidth}
    \textbf{Columna izquierda}
    \begin{itemize}
      \item Elemento 1
      \item Elemento 2
    \end{itemize}
    \column{0.5\textwidth}
    \textbf{Columna derecha}
    \begin{itemize}
      \item Elemento 1
      \item Elemento 2
    \end{itemize}
  \end{columns}
\end{frame}

\section{Conclusión}
{
\setbeamertemplate{background}{\begin{tikzpicture}[remember picture,overlay]
  \fill[verde1] (current page.south west) rectangle (current page.north east);
  \fill[verde2] ([xshift=5cm]current page.north) -- (current page.north east) -- (current page.south east) -- ([xshift=1cm]current page.south) -- cycle;
  \fill[verde3] ([xshift=7cm]current page.north) -- (current page.north east) -- (current page.south east) -- ([xshift=4cm]current page.south) -- cycle;
\end{tikzpicture}}
\begin{frame}[plain]
\begin{tikzpicture}[remember picture,overlay,line join=round,line cap=round]
  \begin{scope}[shift={([xshift=-4.2cm,yshift=-2.2cm]current page.center)},scale=0.9]
    % sombra
    \fill[verdeosc,opacity=0.25] (0.9,-0.15) ellipse (2.4 and 0.22);
    \tikzset{piel/.style={fill=verde5,draw=verdeosc,very thick}}
    % cola
    \fill[piel] (-0.2,2.1) .. controls (-1.2,1.9) and (-2.2,1.2) .. (-2.9,0.6)
      .. controls (-2.0,0.9) and (-1.0,1.2) .. (0.0,1.3) -- cycle;
    % piernas (antes del cuerpo, para que su parte de arriba quede tapada)
    \fill[piel] (0.4,1.3) -- (0.35,-0.1) -- (0.95,-0.1) -- (0.95,1.3) -- cycle;
    \fill[piel] (1.15,1.3) -- (1.1,-0.1) -- (1.7,-0.1) -- (1.65,1.3) -- cycle;
    % cuerpo inclinado
    \fill[piel,rotate around={30:(0.8,1.9)}] (0.8,1.9) ellipse (1.35 and 0.85);
    % cuello
    \fill[piel] (1.2,2.3) -- (1.7,3.4) -- (2.5,3.1) -- (1.9,1.9) -- cycle;
    % muslo (dentro del cuerpo, sobre la pierna)
    \fill[piel] (0.75,1.55) circle (0.48);
    % garras de los pies
    \foreach \x in {0.37,0.6,0.83,1.12,1.35,1.58} \fill[white,draw=verdeosc,thin] (\x,-0.1) -- ++(0.06,-0.12) -- ++(0.06,0.12) -- cycle;
    % mandíbula inferior (boca abierta)
    \fill[piel,rounded corners=6pt] (1.9,3.2) -- (3.5,3.05) -- (3.35,2.7) -- (2.0,2.7) -- cycle;
    \fill[pink!60,rounded corners=4pt] (2.3,3.15) -- (3.35,3.06) -- (3.2,2.9) -- (2.4,2.92) -- cycle;
    % dientes inferiores
    \foreach \x in {2.55,2.85,3.15} \fill[white,draw=verdeosc,thin] (\x,3.07) -- ++(0.07,0.15) -- ++(0.07,-0.15) -- cycle;
    % cráneo
    \fill[piel,rounded corners=10pt] (1.6,3.2) -- (1.75,4.1) -- (3.75,4.0) -- (3.9,3.3) -- cycle;
    % dientes superiores
    \foreach \x in {2.4,2.7,3.0,3.3,3.6} \fill[white,draw=verdeosc,thin] (\x,3.27) -- ++(0.07,-0.17) -- ++(0.07,0.17) -- cycle;
    % ojo y fosa nasal
    \fill[white,draw=verdeosc,thick] (2.45,3.72) circle (0.2);
    \fill[black] (2.52,3.73) circle (0.1); \fill[white] (2.55,3.77) circle (0.03);
    \fill[verdeosc] (3.6,3.82) ellipse (0.07 and 0.04);
    % bracitos
    \draw[verdeosc,line width=4pt] (1.95,2.35) -- (2.35,2.15) -- (2.45,2.3);
    \draw[verde5,line width=2pt] (1.95,2.35) -- (2.35,2.15) -- (2.45,2.3);
    % manchas
    \fill[verde3] (0.4,2.1) circle (0.2); \fill[verde3] (1.0,2.4) circle (0.14);
    \fill[verde3] (0.9,1.7) circle (0.17); \fill[verde3] (2.6,3.55) circle (0.08);
    % púas en la espalda
    \foreach \x/\y/\r in {-0.6/1.85/20,0.0/2.25/25,0.55/2.6/30,1.05/2.95/40,1.45/3.45/55}
      \fill[verdeosc,rotate around={\r:(\x,\y)}] (\x,\y) -- ++(0.14,0.26) -- ++(0.14,-0.26) -- cycle;
    % globo de diálogo: primero la punta, luego el globo encima
    \fill[white,draw=verdeosc,very thick] (6.4,4.6) -- (6.4,4.0) -- (4.0,3.4) -- cycle;
    \node[fill=white,draw=verdeosc,very thick,rounded corners=12pt,inner sep=10pt,
          font=\Huge\bfseries,text=verdeosc] at (8.2,4.4) {¡Gracias!};
  \end{scope}
\end{tikzpicture}
\end{frame}
}

\end{document}
